\section{Operaciones semánticas en el espacio latente}

\subsection{Codificación y decodificación deterministas}

La mapeo biyectivo $\mathbf{x}_T \leftrightarrow \mathbf{x}_0$ establecido por la muestreo determinista de DDIM permite no solo generar $\mathbf{x}_0$ a partir de $\mathbf{x}_T$ (decodificar), sino también recuperar $\mathbf{x}_T$ a partir de $\mathbf{x}_0$ (codificar). El proceso de codificación se realiza ejecutando la dirección hacia adelante de DDIM:

$$
\mathbf{x}_{t+1} = \sqrt{\bar{\alpha}_{t+1}}\,\hat{\mathbf{x}}_0(\mathbf{x}_t) + \sqrt{1-\bar{\alpha}_{t+1}}\cdot\frac{\mathbf{x}_t - \sqrt{\bar{\alpha}_t}\,\hat{\mathbf{x}}_0(\mathbf{x}_t)}{\sqrt{1-\bar{\alpha}_t}}
$$

Esto significa que, para cualquier imagen real $\mathbf{x}_0$, podemos encontrar su ``punto correspondiente'' $\mathbf{x}_T$ en el espacio de ruido; ejecutando la decodificación DDIM a partir de dicho punto, se puede recuperar la imagen original con precisión.

\begin{importantbox}{Estructura semántica del espacio latente de DDIM}
Debido a la existencia del mapeo determinista, el espacio de ruido $\mathbf{x}_T$ se convierte efectivamente en un espacio latente con estructura semántica. En este espacio:
\begin{itemize}
    \item Representaciones similares de imágenes también son cercanas en el espacio latente
    \item Las operaciones lineales en el espacio latente (como la interpolación) corresponden a cambios semánticos significativos en el espacio de datos
    \item A diferencia del espacio latente de VAE (de baja dimensión), el espacio latente de DDIM tiene la misma dimensión que el espacio de datos
\end{itemize}
\end{importantbox}

\subsection{Interpolación en el espacio latente}

El artículo de DDIM muestra una aplicación interesante de la interpolación en el espacio latente. Dadas dos imágenes reales $\mathbf{x}_0^{(1)}$ y $\mathbf{x}_0^{(2)}$:

\begin{enumerate}
    \item Se obtienen $\mathbf{x}_T^{(1)}$ y $\mathbf{x}_T^{(2)}$ mediante codificación DDIM respectivamente
    \item Se realiza una interpolación lineal esférica (SLERP) en el espacio latente:
    $$
    \mathbf{x}_T^{(\lambda)} = \frac{\sin((1-\lambda)\theta)}{\sin\theta}\,\mathbf{x}_T^{(1)} + \frac{\sin(\lambda\theta)}{\sin\theta}\,\mathbf{x}_T^{(2)}, \quad \lambda \in [0, 1]
    $$
    donde $\theta = \arccos\left(\frac{\langle\mathbf{x}_T^{(1)}, \mathbf{x}_T^{(2)}\rangle}{\|\mathbf{x}_T^{(1)}\|\,\|\mathbf{x}_T^{(2)}\|}\right)$
    \item Se ejecuta la decodificación DDIM en el punto interpolado $\mathbf{x}_T^{(\lambda)}$ para obtener la imagen interpolada
\end{enumerate}

\begin{knowledgebox}{¿Por qué se usa interpolación esférica (SLERP) en lugar de lineal?}
Dado que los vectores de ruido $\mathbf{x}_T$ en el espacio latente se muestrean de una distribución gaussiana estándar, en espacios de alta dimensión estos vectores se distribuyen aproximadamente uniformemente sobre una hiperesfera (la norma se concentra cerca de $\sqrt{d}$). La interpolación lineal (LERP) genera puntos intermedios con normas menores, desviándose del conjunto típico de la distribución gaussiana y lo que podría resultar en una disminución de la calidad. La interpolación esférica (SLERP) mantiene constante la norma del vector; los puntos intermedios generados permanecen en la hiperesfera, por lo que es más razonable.
\end{knowledgebox}

\subsection{Comparación con otros espacios latentes}

\begin{table}[H]
\centering
\caption{Comparativa de características del espacio latente entre modelos diferentes}
\begin{tabular}{lccc}
\toprule
\textbf{Modelo} & \textbf{Dimensión del espacio latente} & \textbf{Codificación determinista} & \textbf{Calidad de interpolación} \\
\midrule
VAE & Baja dimensión & Aleatorio & Medio \\
GAN & Baja dimensión & Generalmente sin codificador & Buena \\
DDIM & Igual que los datos & Determinista & Buena \\
\bottomrule
\end{tabular}
\end{table}

\begin{warningbox}{La precisión de la codificación DDIM depende del número de pasos}
El proceso de codificación DDIM es esencialmente una aproximación numérica de la integración de ODE. Cuantos menos pasos se utilicen, mayor será el error numérico y menor la calidad de reconstrucción del ciclo ``codificación $\to$ decodificación''. En aplicaciones que requieren reconstrucción precisa (como la edición de imágenes), generalmente es necesario utilizar un mayor número de pasos (por ejemplo, entre 100 y 200) para garantizar la precisión de la codificación.
\end{warningbox}

\subsection{Resumen de la sección}

El muestreo determinista de DDIM construye un mapeo biyectivo desde el espacio de ruido hacia el espacio de datos, dotando al espacio de ruido de una estructura semántica. Mediante el flujo de codificación-interpolación-decodificación, es posible lograr una interpolación de imágenes de alta calidad. La interpolación lineal esférica (SLERP) resulta más adecuada para espacios latentes con distribución gaussiana que la interpolación lineal.
